% Generated draft: requires mathematical and source review.
\documentclass[11pt]{article}
\usepackage[margin=25mm]{geometry}
\usepackage{fontspec}
\setmainfont{Latin Modern Roman}
\usepackage{amsmath,amssymb,mathtools,unicode-math}
\setmathfont{Latin Modern Math}
\usepackage{enumitem,needspace,xurl,hyperref}
\setlength{\parindent}{0pt}
\setlength{\parskip}{5pt}
\setlength{\emergencystretch}{3em}
\newcommand{\N}{\mathbb{N}}
\newcommand{\Z}{\mathbb{Z}}
\newcommand{\Q}{\mathbb{Q}}
\newcommand{\R}{\mathbb{R}}
\newcommand{\C}{\mathbb{C}}
\newcommand{\F}{\mathbb{F}}
\newcommand{\A}{\mathbb{A}}
\newcommand{\PP}{\mathbb{P}}
\newcommand{\E}{\mathbb{E}}
\newcommand{\eps}{\varepsilon}
\newcommand{\dd}{\,\mathrm{d}}
\begin{document}
\section*{UnsolvedMath ID: 30004198}
\textbf{Source problem code:} catalogue entry
\subsection{Definitions and mathematical statement}
Let \(\mathcal C=\mathrm{SH}(\C)^{\wedge}_2\) denote the \(2\)-complete motivic stable homotopy category over \(\C\), interpreted with its symmetric monoidal stable \(\infty\)-categorical enhancement when discussing commutative algebras. Write \(S^0\) for the motivic sphere and
\[
C_\tau=S^0/\tau=\operatorname{cofib}(\tau\colon A\longrightarrow S^0)
\]
for the motivic cofiber intended by the source. Here \(A\) denotes the domain of the intended motivic map \(\tau\); the supplied record does not specify its degree or construction. Regard \(C_\tau\) as the commutative algebra specified in the source, hence as an object of \(\operatorname{CAlg}(\mathcal C)\).

A normed algebra carries compatible multiplicative norm maps along the finite \'{e}tale morphisms in its indexing category. In such a framework, for an allowed finite \'{e}tale morphism \(f\colon T\to X\), these include maps \(f_\otimes E_T\to E_X\), where \(f_\otimes\) is the motivic norm functor, with coherent compatibility with composition, base change, and the algebra structures. The supplied record does not specify the indexing category or how norms are to interact with \(2\)-completion.

\textbf{Question.} In the intended normed-algebra framework over \(\C\), let \(\mathcal N\) be its category of normed algebras and let
\[
U\colon\mathcal N\longrightarrow\operatorname{CAlg}(\mathcal C)
\]
be the functor forgetting the norm data. Does there exist \(\widetilde C_\tau\in\mathcal N\) together with an equivalence
\[
U(\widetilde C_\tau)\simeq C_\tau
\qquad\text{in }\operatorname{CAlg}(\mathcal C)?
\]
This expresses the lifting question while retaining the underlying commutative algebra structure. A fully specified formulation requires identification of \(\tau\), the intended normed-algebra category, and its relationship to \(2\)-completion; these choices are unresolved in the supplied record.
\subsection{Short English statement}
Does the commutative algebra structure on the \(2\)-complete motivic cofiber \(S^0/\tau\) over \(\C\) extend to a compatible normed algebra structure? The supplied record leaves the precise definition of \(\tau\) and the normed framework unspecified.
\subsection{Sources}
\begin{itemize}
\item Supplied local UnsolvedMath registry record, permanent ID \(30004198\), title \emph{Normed Algebra Lifts of the Motivic Tau Cofiber}, source problem code \texttt{catalogue entry}, in the Oberwolfach Reports Open Problems collection. The record gives proposed year \(2019\) and identifies no proposer.
\item Supplied report citation: \emph{Homotopy Theory (2020)}, Oberwolfach Reports, pp.\ 2183--2256, DOI \texttt{10.4171/owr/2019/36}. Source link: \url{https://doi.org/10.4171/owr/2019/36}. The background attributes the problem to workshop open-problem discussions documented in this report.
\item Supplied literature source, labelled ``original Oberwolfach report'': \url{https://doi.org/10.4171/OWR/2019/36}.
\item Supplied literature source, \emph{The Motivic Cofiber of tau}: \url{https://doi.org/10.4171/DM/642}.
\item Literature assessment dated exactly \texttt{2026-08-22}: the supplied curation labels the problem open and states that the motivic cofiber has a unique \(E_\infty\) ring structure, while a lift to a normed-algebra category remains open. This is supplied dated curation derived from the OWR record, not an independently verified literature assessment. The record also labels the statement \texttt{corrected\_verified}, explaining that notation, typography, and sentence structure were repaired while preserving mathematical content; this label is not independent verification of the present draft.
\end{itemize}
\end{document}
